\subsection{Sharp bounds for a scalar mean chain}\label{sec:wholechain}

A whole-chain factory improves on independent layer composition
when the network depends on the input through one scalar mean.
This section proves matching query bounds for that family.

Let \(F_{s,0}(z)=z\), \(F_{s,d+1}(z)=\tanh(sF_{s,d}(z))\),
where \(s>0\) is dyadic. The analytic identities below hold for
every positive real \(s\); the dyadic assumption supplies an
executable parameter description.
A source request returns an independent sign of
unknown mean \(z\in[-1,1]\).
Let \(U_D(s)\) be the infimum, over exact factories for a sign of
mean \(F_{s,D}(z)\), of the worst-mean expected source count.
Auxiliary randomness and arithmetic are excluded from this
particular query measure.

\begin{theorem}[Uniform scalar complexity]\label{thm:uniformchain}
Set \(\lambda=s^D\) and \(A=\sum_{j=1}^D s^{2j}\).
Then, with universal constants,
\begin{equation}\label{eq:uniformchain}
\frac1{25}\lambda\sqrt{1+A}
\le U_D(s)\le2\lambda\sqrt{1+A}.
\end{equation}
Thus the dependence on the gain and depth is determined jointly,
including when \(s\) approaches one with \(D\).
\end{theorem}

\begin{theorem}[The scalar phase diagram]\label{thm:chainphase}
For \(D\ge1\),
\begin{align}
(1-s)s^D\le U_D(s)&\le\frac{s^D}{1-s^2},
&&0<s<1,\label{eq:chainsub}\\
\frac{\sqrt D}{8}\le U_D(1)&\le2\sqrt D,\label{eq:chaincritical}\\
s^{2D}\le U_D(s)&\le
\frac{2s}{\sqrt{s^2-1}}s^{2D},
&&s>1.\label{eq:chainsuper}
\end{align}
The upper bounds are realized by an exact majority mixture for
the entire composition.
\end{theorem}

\subsubsection{Positivity for the entire composition}

A smooth function \(f\) on \((0,\infty)\) is
\emph{completely monotone} if \((-1)^kf^{(k)}\ge0\) for all \(k\).
A nonnegative smooth function is a \emph{Bernstein function} if
its derivative is completely monotone.
Two elementary closure facts will suffice. Products of completely
monotone functions are completely monotone, by Leibniz's rule.
If \(C\) is completely monotone and \(g\) is Bernstein, then
\(C\circ g\) is completely monotone. Indeed, a term in the
Faà di Bruno formula has the form
\[
C^{(k)}(g)\prod_{j\ge1}(g^{(j)})^{m_j},
\quad \sum_jm_j=k,\quad\sum_jjm_j=q,
\]
and its sign is \((-1)^{k+\sum_j(j-1)m_j}=(-1)^q\).
The coefficients in that formula are nonnegative.
Differentiating a composition then shows that Bernstein functions
are closed under composition.

Put
\[
C(w)=\sech^2\sqrt w,\qquad
T(w)=\frac{\tanh\sqrt w}{\sqrt w},\qquad
g(w)=\tanh^2\sqrt w.
\]
The cosh product shows that \(C\) is completely monotone:
each finite product of factors \((1+\alpha w)^{-2}\) has this
property, and local uniform convergence permits differentiation.
Its logarithmic derivative gives
\[
T(w)=\sum_{j\ge0}\frac8{\pi^2(2j+1)^2+4w},
\]
which is completely monotone as well.
Since \(g'=TC\), \(g\) is Bernstein.
Positive scaling preserves these properties.
Writing \(g_s(w)=g(s^2w)\), \(C_s(w)=C(s^2w)\), we obtain
\begin{equation}\label{eq:chainH}
H_{D,s}(w):=F_{s,D}'(\sqrt w)
=s^D\prod_{j=0}^{D-1}C_s(g_s^{\circ j}(w)),
\end{equation}
a completely monotone function.

This derivative has the analytic domain needed for its Taylor
series. If \(\sqrt w=x+iy\) with \(\Re w>0\), then \(x>|y|\), and
\[
\Re\tanh^2(x+iy)
=\frac{\sinh^2(2x)-\sin^2(2y)}
       {(\cosh(2x)+\cos(2y))^2}>0.
\]
Thus \(g_s\) maps the right half-plane into itself, and
\(H_{D,s}\) is holomorphic there.
Its removable value at zero is \(s^D\).
Consequently
\[
H_{D,s}(1-v)=\sum_{k\ge0}d_kv^k,\qquad d_k\ge0,\qquad
\sum_kd_k=s^D.
\]
The coefficient signs follow from complete monotonicity at \(w=1\);
the sum follows by monotone convergence as \(v\uparrow1\).

Let \(R=F_{s,D}(1)>0\).
The same derivative integration as in Lemma~\ref{lem:majority} proves
\begin{equation}\label{eq:chainmixture}
\frac{F_{s,D}(z)}R=\sum_{k\ge0}\pi_kM_k(z),
\qquad \pi_k=\frac{d_k}{R C_k},\qquad \sum_k\pi_k=1.
\end{equation}
Open a gate of probability \(R\). On a closed gate return a fair
sign; on an open gate draw \(K\sim\pi\) and return its majority.
This sign has mean \(F_{s,D}(z)\).

\subsubsection{The cost of the global mixture}

Define \(A_D(s)=\sum_{j=1}^D s^{2j}\).
We count every vote of the full majority; stopping early can only
lower the cost. Integration by parts gives the exact full-majority mean
\begin{equation}\label{eq:chainarity}
Q_{D,s}=s^D+
\int_0^1\frac{s^D-F_{s,D}'(z)}{z^2}\,dz.
\end{equation}
The gate probability \(R\) cancels the normalization in
\eqref{eq:chainmixture}.
For \(z\ge0\), \(F_{s,j}(z)\le s^jz\), and
\[
\frac{F_{s,D}'(z)}{s^D}
=\prod_{j=1}^D(1-F_{s,j}(z)^2).
\]
Using \(1-\prod_j(1-u_j)\le\sum_ju_j\) for \(0\le u_j\le1\),
\[
0\le s^D-F_{s,D}'(z)
\le s^D\min\{1,A_D(s)z^2\}.
\]
Hence
\begin{equation}\label{eq:chainQ}
Q_{D,s}\le s^D(1+A_D(s)),\qquad
Q_{D,s}\le2s^D\sqrt{A_D(s)}\quad\text{if }A_D(s)\ge1.
\end{equation}
The first bound holds for every \(A_D(s)\).
It implies the upper bound in \eqref{eq:uniformchain}: if \(A\le1\),
\(1+A\le2\sqrt{1+A}\), and the second bound handles \(A\ge1\).
For \(s<1\), \(1+A_D(s)\le(1-s^2)^{-1}\).
At \(s=1\), \(A_D(1)=D\).
For \(s>1\),
\(A_D(s)\le s^2s^{2D}/(s^2-1)\).
These inequalities prove all upper bounds in
Theorem~\ref{thm:chainphase}.

\subsubsection{Lower bounds without a stopping-time regularity assumption}

Consider any almost surely terminating scalar factory whose
worst-mean expected source count is \(M<\infty\).
Truncate it after \(T\) source requests and return a fair sign if
it has not terminated. Its mean \(f_T\) is a polynomial and
\[
\sup_{z\in[-1,1]}|f_T(z)-f(z)|\le M/T.
\]
The finite transcript proof of Lemma~\ref{lem:transcript} gives
\[
|f_T'(z)|\le\sqrt{\frac{M}{1-z^2}},\qquad
|f_T''(z)|\le\frac{2M}{1-z^2}.
\]
Integrate the first bound on a short interval and pass to the
uniform limit. For the second, use
\[
f_T(z+h)-2f_T(z)+f_T(z-h)
=\int_{-h}^h(h-|u|)f_T''(z+u)\,du.
\]
Pass to the limit in \(T\), divide by \(h^2\), and then let
\(h\to0\). For twice differentiable \(f\),
\begin{equation}\label{eq:unknownlower}
M\ge(1-z^2)f'(z)^2,\qquad
M\ge\frac{1-z^2}{2}|f''(z)|.
\end{equation}
These are bounds on the global worst-mean cost \(M\).
No differentiation through an unbounded runtime is needed.

\paragraph{Proof idea.}
The reciprocal-square recurrence compares the chain's range with its
derivative at zero. An input interval on which the derivative stays large
then forces any exact factory to use enough source samples.

\begin{proof}[Lower bound in Theorem~\ref{thm:uniformchain}]
The reciprocal-square inequalities
\[
u^{-2}+\frac23\le\coth^2u\le u^{-2}+1,
\]
from \eqref{eq:cothradius} and \(\sinh u\ge u\), give the range
comparison
\begin{equation}\label{eq:chainrange}
\frac{\lambda}{\sqrt{1+A}}\le R=F_{s,D}(1)
\le\frac{\lambda}{\sqrt{1+\frac23A}}.
\end{equation}
Indeed, apply the inequalities to each step of the radius recurrence
and multiply the final reciprocal inequality by \(s^{2D}\).
The accumulated geometric sum is exactly \(A\).

Endpoint coupling (Section~\ref{sec:lower}) gives \(M\ge R\).
If \(A<24\), this is at least
\(\lambda\sqrt{1+A}/25\).
If \(A\ge24\), equation \eqref{eq:chainrange} gives
\(T=2R/\lambda\le2/\sqrt{17}<1/2\). The target has derivative
\(\lambda\) at zero and range \([-R,R]\), so the bending
bound \eqref{eq:bending}, with the cost bound
\eqref{eq:unknownlower}, gives
\[
M\ge\frac{3\lambda^2}{16R}
\ge\frac3{16}\sqrt{\frac23}\lambda\sqrt{1+A}
\ge\frac1{25}\lambda\sqrt{1+A}.
\]
\end{proof}

For \(f=F_{s,D}\), \(f'(0)=s^D\), proving the supercritical lower
bound. For \(s=1\), \eqref{eq:tanhlower} gives
\(F_{1,j}(z)^2\ge(z^{-2}+j)^{-1}\).
Take \(z_0=1/(2\sqrt D)\). Then
\[
F_{1,D}'(z_0)
=\prod_{j=1}^D(1-F_{1,j}(z_0)^2)
\le(1-1/(5D))^D\le e^{-1/5}\le5/6.
\]
Since \(F_{1,D}'(0)=1\), the mean value theorem gives some
\(z_*\in(0,z_0)\) with
\(|F_{1,D}''(z_*)|\ge\sqrt D/3\).
Here \(1-z_*^2\ge3/4\), so \eqref{eq:unknownlower} yields
\(M\ge\sqrt D/8\).

For \(0<s<1\), endpoint coupling with all-plus and all-minus
source signs shows that some source request occurs with probability
at least \(R=F_{s,D}(1)\), the difference of the two target
Bernoulli probabilities. For \(u\ge0\), the bound
\(\tanh u\ge u/(1+u)\) follows from \(e^{2u}\ge1+2u\).
Writing \(R_j=F_{s,j}(1)\),
\[
\frac{s^{j+1}}{R_{j+1}}
\le\frac{s^j}{R_j}+s^{j+1}.
\]
Thus \(s^D/R\le1+\sum_{j=1}^D s^j\le(1-s)^{-1}\).
This proves the remaining lower bound.

\subsubsection{Finite inputs and arbitrary depth}

For a vector of \(m\) signs, independently choosing a uniform
coordinate on each request supplies the unknown-mean source with
\(z=m^{-1}\sum_jx_j\). Repeated coordinate choices remain independent
conditional on the fixed input, because the indices are independent.
A second option opens the same gate \(R\), then reads all \(m\)
signs and evaluates the scalar function if the gate opens.
It uses \(mR\) expected probes.

\begin{theorem}[Uniform finite-input complexity]\label{thm:finitejoint}
Let \(m\ge2\) be even, \(D\ge1\), \(s>0\) dyadic,
\(\lambda=s^D\), \(A=\sum_{j=1}^D s^{2j}\), and
\[
\Xi_{s,D,m}=\lambda\min\left\{\sqrt{1+A},
                      \frac{m}{\sqrt{1+A}}\right\}.
\]
For the mean-chain target on \(m\) input signs, the optimal
worst-input expected number \(Q^*_{s,D,m}\) of distinct probes obeys
\begin{equation}\label{eq:finitejoint}
\frac1{128}\Xi_{s,D,m}\le Q^*_{s,D,m}\le2\Xi_{s,D,m}.
\end{equation}
The constants are independent of \(s,D,m\).
\end{theorem}
\paragraph{Proof idea.}
Choose between a global majority factory and a rare complete input read. For
the lower bound, apply the transcript score to the empirical mean and use
its fourth moment to retain the dependence on the finite input size.

\begin{proof}
Choose the global factory when \(1+A\le m\), and the gated scan
otherwise. Equations \eqref{eq:chainQ} and \eqref{eq:chainrange}
give the upper bound in both cases.
This choice uses an exact rational comparison before seeing the input.

For the lower bound, endpoint coupling gives \(Q^*\ge R\).
If \(A<100\), then
\[
R\ge\frac{\lambda}{\sqrt{1+A}}
\ge\frac{\Xi_{s,D,m}}{101}.
\]
Suppose now that \(A,m\ge100\).
Equation \eqref{eq:derivativelower} gives
\(h=H'(0)\ge\lambda/\sqrt{1+3A/m}\).
Using \eqref{eq:chainrange}, the interval \(T=2R/h\) has
\[
T^2\le\frac{4(1+3A/m)}{1+\frac23A}
\le6/A+18/m\le0.24<1/4.
\]
The bending bound \eqref{eq:bending} and the second transcript
bound yield
\[
Q^*\ge\frac{3h^2}{16R}
\ge\frac{3\lambda\sqrt{1+\frac23A}}{16(1+3A/m)}
\ge\frac{3}{64}\sqrt{\frac23}\,\Xi_{s,D,m}.
\]
For the last inequality,
\(\sqrt{1+2A/3}\ge\sqrt{2/3}\sqrt{1+A}\) and
\(1+3A/m\le4\max\{1,(1+A)/m\}\).

It remains to consider \(m<100\).
At a balanced input, flipping one coordinate changes the plus
probability by \(F_{s,D}(2/m)/2\). Lemma~\ref{lem:attentioncoupling},
summed over the \(m\) coordinates, gives
\[
Q^*\ge\frac m2F_{s,D}(2/m)
\ge\frac{\lambda}{\sqrt{1+4A/m^2}}.
\]
If \(A\le m^2\), this is at least \(\lambda/\sqrt5\), while
\(\Xi_{s,D,m}\le\lambda\sqrt m<10\lambda\).
If \(A>m^2\), it is at least
\(m\lambda/\sqrt{5A}\ge\Xi_{s,D,m}/\sqrt5\).
Each displayed constant is larger than \(1/128\).
\end{proof}

\begin{theorem}[Finite critical mean chain]\label{thm:finitecritical}
For even \(m\ge2\), \(D\ge1\), and
\(f(x)=F_{1,D}(m^{-1}\sum_jx_j)\), the optimal worst-input
expected number of distinct probes lies between
\[
\frac1{100}\min\{\sqrt D,m/\sqrt D\}
\quad\text{and}\quad
2\min\{\sqrt D,m/\sqrt D\}.
\]
For every fixed \(s>1\), the corresponding order is
\(\Theta_s(\min\{s^{2D},m\})\).
\end{theorem}
\paragraph{Proof idea.}
Combine the two available upper bounds. For the lower bound, run the
case split of the preceding proof at \(s=1\) and track the constants
against \(\sqrt D\), including the regime where the number of inputs
is smaller than the depth.

\begin{proof}
The critical upper bound follows by choosing between
\(Q_{D,1}\le2\sqrt D\) and \(mR_D\le2m/\sqrt D\).
The supercritical claim follows from
\eqref{eq:lowersuper}, \eqref{eq:chainQ}, and \(mR\le m\).

For the critical lower bound set
\(\psi=\min\{\sqrt D,m/\sqrt D\}\), and use the proof of
Theorem~\ref{thm:finitejoint} at \(s=1\), where \(\lambda=1\) and
\(A=D\). If \(D<100\), endpoint coupling and
\eqref{eq:chainrange} give
\(\sup_x\E[Q\mid x]\ge R_D\ge(1+D)^{-1/2}\ge1/10\),
while \(\psi\le\sqrt D<10\). If \(D,m\ge100\), its bending case gives
\[
\sup_x\E[Q\mid x]\ge\frac{3\sqrt{1+2D/3}}{16(1+3D/m)}
\ge\frac{3}{64\sqrt{3/2}}\,\psi>\frac{\psi}{100},
\]
since \(\sqrt{1+2D/3}\ge\sqrt{2/3}\sqrt D\) and
\(\sqrt D/(1+3D/m)\ge\psi/4\).
If \(m<100\), its balanced-input bound reads
\begin{equation}\label{eq:balancedinfluence}
\E[Q\mid x_{\rm balanced}]
\ge\frac m2F_{1,D}(2/m)
\ge(1+4D/m^2)^{-1/2}.
\end{equation}
If \(D\le m^2\), this is at least \(1/\sqrt5\), while
\(\psi\le\sqrt m<10\).
If \(D>m^2\), it is at least
\(m/\sqrt{5D}=\psi/\sqrt5\).
All cases imply the stated constant.
\end{proof}

For arbitrary critical networks one may also enter the root factory
with probability \(R\le r_D\), then read the whole context and
evaluate its conditional scalar probability. This gives at most
\(2n/\sqrt{D+1}\) expected input probes.
Together with Theorem~\ref{thm:finitecritical}, it proves full-class
query order \(\Theta(n/\sqrt D)\) for \(D\ge n\), whenever the
uniform first row is on the parameter grid.
This statement concerns queries; evaluating on the rare open branch
may still be arithmetically expensive.
